% ==============================================================================
% Fichier     : 05_types_audits.tex
% Description : Principaux Types d'Audit dans les Systèmes d'Information
% ==============================================================================

\chapter{Principaux Types d'Audit dans les Systèmes d'Information}%
\label{chap:05_types_audits}

\begin{boxobjectifs}
À l'issue de ce chapitre, l'apprenant doit être capable de~:
\begin{itemize}
    \item Distinguer les huit types d'audit SI dans leurs périmètres, enjeux
        et déclencheurs respectifs~;
    \item Identifier les référentiels internationaux associés à chaque type~;
    \item Comprendre les interactions et dépendances entre types d'audits~;
    \item Positionner chaque type dans le continuum
        \textit{gouvernance, fonctionnement, sécurité, résilience}.
\end{itemize}
\end{boxobjectifs}

L'audit des systèmes d'information n'est pas un acte monolithique. Il recouvre un
spectre de missions dont la nature, le périmètre, les enjeux et la posture de
l'auditeur diffèrent profondément selon l'objet examiné. Auditer une application
comptable n'appelle pas les mêmes compétences ni les mêmes outils qu'auditer une
infrastructure réseau ou un projet en cours de développement. Chaque type d'audit
répond à une question fondamentale distincte~: \textit{le logiciel fait-il ce qu'on
attend de lui~?}, \textit{la DSI est-elle gouvernée efficacement~?}, \textit{le projet
sera-t-il livré conformément aux engagements~?}, \textit{le SI résistera-t-il à une
cyberattaque ou à un sinistre~?}

Ce chapitre présente les huit grands types d'audit SI rencontrés en milieu
professionnel. Pour chacun, la même structure est suivie~: enjeu fondamental,
périmètre et frontières, fréquence et déclencheurs, dimensions caractéristiques,
référentiels mobilisés, interactions avec les autres types et illustration par le
cas LiZbiZ. Les décompositions opérationnelles détaillées, points de contrôle,
critères, questions d'évaluation, risques et conséquences, font l'objet du chapitre
suivant.

\begin{boxinfo}
\textbf{Une lecture systémique.} Les types d'audit ne sont pas étanches. Un audit
de la fonction informatique révèle souvent des lacunes dans la gestion des projets~;
un audit de sécurité met fréquemment en lumière des déficiences organisationnelles
que seul un audit de la DSI aurait pu formellement documenter. L'auditeur SI
expérimenté perçoit ces connexions et sait élargir son périmètre lorsque les
constats le justifient.
\end{boxinfo}


% ==============================================================================
\section{Audit des Applications en Service}
\label{sec:audit-appli}
% ==============================================================================

\subsection{Enjeu fondamental}

Une application en service est un logiciel opérationnel accompagnant, automatisant
ou se substituant à un processus de l'organisation, qu'il s'agisse d'un ERP gérant
les flux financiers, d'un CRM pilotant la relation client, d'un progiciel de paie ou
d'un système de gestion des stocks. Ces applications sont au cœur des processus
métiers~: leurs défaillances se traduisent directement en pertes financières, en
erreurs de traitement ou en violations réglementaires.

L'enjeu de cet audit est double. D'un côté, il s'agit d'émettre une
\termecle{assurance raisonnable} sur la fiabilité de l'application, c'est-à-dire
sur la qualité de son contrôle interne et la validité des données qu'elle traite
et restitue. De l'autre, il s'agit d'apprécier son adéquation aux besoins de
l'organisation, sa performance et sa capacité à évoluer avec elle. Ces deux
périmètres, fiabilité-sécurité d'une part, efficacité-performance de l'autre,
constituent les deux visées distinctes mais complémentaires de l'audit d'application
en service. Un audit complet les couvre toutes deux.

\subsection{Périmètre et frontières}

Le périmètre de cet audit couvre l'application elle-même dans toutes ses
composantes~: programmes, données, paramètres, documentation et habilitations.
Il englobe également l'environnement dans lequel elle s'exécute, système
d'exploitation, base de données, infrastructure réseau, et les processus
organisationnels qui entourent son exploitation~: gestion des accès, procédures
de sauvegarde, plans de continuité, contrats de maintenance.

La frontière avec l'audit de la fonction informatique mérite d'être précisée.
L'audit d'application se concentre sur \textit{ce que fait} le logiciel et
\textit{comment il est exploité}~; l'audit de la DSI examine \textit{comment est
gouvernée} l'entité qui le produit et le maintient. En pratique, les deux missions
se nourrissent mutuellement~: une faiblesse détectée dans la gestion des droits
d'accès d'une application peut révéler une déficience structurelle dans la politique
de sécurité de la DSI.

\subsection{Dimensions caractéristiques}

L'audit d'application en service s'articule autour de deux axes analytiques
fondamentaux.

\textbf{Les contrôles généraux} (\textit{General Controls}) couvrent l'environnement
global dans lequel l'application fonctionne~: sécurité du système d'exploitation hôte,
protection physique des serveurs, procédures de sauvegarde et de reprise, gestion
des mises à jour et des correctifs, séparation entre les environnements de
développement, de test et de production. Ces contrôles conditionnent la fiabilité
de tous les contrôles applicatifs~: si l'environnement est compromis, aucun contrôle
programmatique ne peut être considéré comme fiable.

\textbf{Les contrôles d'application} (\textit{Application Controls}) portent sur
le logiciel métier lui-même et s'organisent selon le modèle \termecle{IPO — Input,
Process, Output}. En entrée, l'auditeur vérifie que seules des données complètes,
exactes et autorisées sont acceptées par le système. En traitement, il s'assure que
les calculs et règles de gestion sont correctement implémentés et que les données
sensibles sont journalisées. En sortie, il contrôle que les résultats produits sont
distribués aux seuls destinataires habilités et que les états de sortie font l'objet
d'un contrôle utilisateur effectif.

La \termecle{piste d'audit} (\textit{audit trail}) constitue l'épine dorsale de tout
audit d'application~: toute action significative sur les données, création,
modification, suppression, doit être horodatée, attribuée à un utilisateur identifié
et conservée de façon intègre. Son absence est une observation critique majeure.

% ------------------------------------------------------------------------------
\subsection{Au-delà du Modèle IPO~: les Contrôles des Interfaces Inter-applicatives}
\label{subsec:controles-interfaces}
% ------------------------------------------------------------------------------

Le modèle IPO décrit le cycle de traitement \emph{interne} à une application. Mais
un système d'information d'entreprise est rarement composé d'une seule application
autosuffisante~: il est un réseau d'applications qui échangent des données,
automatiquement, semi-automatiquement ou manuellement, via des interfaces. Ces
échanges inter-applicatifs constituent une source de risque propre, distincte des
risques internes à chaque application~: une donnée peut être correctement produite
par l'application source et corrompue, tronquée ou détournée lors de son transfert
vers l'application destination.

L'interface n'est pas un 4\up{ème} niveau séquentiel du modèle IPO~: elle peut se
greffer à n'importe quel point du cycle (alimenter un Input, déclencher un traitement
ou transmettre un Output). Elle constitue un \textbf{domaine de contrôle transversal},
distinct et complémentaire de la matrice IPO.

\begin{boxdef}
\textbf{Interface inter-applicative~:}
Lien d'échange de données entre deux applications ou systèmes distincts, caractérisé
par~: son \emph{type} (automatique, semi-automatique ou manuel)~; les applications
source et destination~; la \emph{nature des flux} échangés~; son \emph{protocole
d'échange} (fichier plat, API REST, Web Service, EDI, message queue)~; sa
\emph{périodicité} (temps réel, quotidien, hebdomadaire)~; et ses \emph{contrôles
intégrés} (totaux de contrôle, accusé de réception, chiffrement). Plus une interface
est manuelle ou semi-automatique, plus son niveau de risque est élevé~: chaque
intervention humaine dans un flux de données est un point potentiel d'erreur,
d'omission ou de fraude.
\end{boxdef}

\subsubsection{Cartographie des Interfaces~: Préalable à l'Audit}

Avant tout test, l'auditeur constitue la cartographie des interfaces de l'application
auditée. Pour chaque interface identifiée, il documente~:
\begin{itemize}
    \item Le type~: \textbf{automatique}, \textbf{semi-automatique} ou
        \textbf{manuel}~;
    \item Les applications source et destination~;
    \item Le protocole, le format d'échange et la périodicité~;
    \item Le responsable désigné pour le traitement des anomalies~;
    \item Les contrôles intégrés~: totaux de contrôle, numérotation de séquence,
        accusé de réception, chiffrement.
\end{itemize}

\paragraph{Signal de risque~: l'interface non inventoriée.}
Les interfaces les plus dangereuses sont celles qui ne figurent pas dans la
cartographie officielle~: transferts de fichiers par clé USB, exports Excel
réintégrés manuellement dans un ERP, flux échangés par messagerie électronique.
Ces \emph{interfaces informelles} échappent à tout contrôle intégré et constituent
des angles morts de l'audit applicatif. L'auditeur doit donc non seulement valider
la cartographie fournie par l'audité, mais interroger les utilisateurs sur leurs
pratiques réelles d'échange de données.

\subsubsection{Les Objectifs de Contrôle des Interfaces}

Trois objectifs structurent l'audit des interfaces.

\paragraph{1. L'intégrité des flux.}
Les données transmises doivent être reçues de façon \emph{exhaustive},
\emph{exacte} et \emph{dans l'ordre attendu}~:
\begin{itemize}
    \item \textbf{Totaux de contrôle de lots (\textit{batch control totals})~:}
        la somme des montants ou le nombre d'enregistrements calculés côté source
        est comparé à la valeur reçue côté destination~; toute discordance déclenche
        une alerte~;
    \item \textbf{Numérotation de séquence~:} chaque enregistrement porte un numéro
        séquentiel~; l'application destination détecte les ruptures (perte) et les
        doublons (duplication)~;
    \item \textbf{Rapprochement périodique~:} réconciliation formelle entre les
        soldes ou totaux des deux systèmes à une fréquence définie et documentée.
\end{itemize}

\paragraph{2. La confidentialité et la non-répudiation des échanges.}
Lorsque les données échangées sont sensibles~:
\begin{itemize}
    \item Les données sont \textbf{chiffrées en transit}~: protocole TLS~1.2 minimum
        pour les échanges réseau~;
    \item \textbf{L'identité de l'émetteur est authentifiée}~: un certificat ou une
        clé partagée garantit que le flux provient bien du système source déclaré~;
    \item Les échanges sont \textbf{non répudiables}~: une signature numérique ou
        un journal horodaté permet de prouver qu'un message a bien été envoyé par
        telle application à telle heure et reçu par telle autre.
\end{itemize}

\paragraph{3. La résilience de l'interface et le traitement des anomalies.}
\begin{itemize}
    \item Un \textbf{responsable d'interface est désigné}~: une personne nommée
        surveille les interfaces et traite les anomalies dans les délais définis~;
    \item Les \textbf{anomalies sont journalisées et traitées}~: un registre des
        erreurs est tenu, les anomalies répétitives signalent un défaut de conception
        à corriger en profondeur~;
    \item Une \textbf{procédure de rejeu} existe~: en cas d'échec d'un transfert,
        la procédure permet de relancer le flux sans créer de doublon ni perdre de
        données~;
    \item Les flux sont \textbf{testés avant mise en production}~: aucune interface
        nouvelle ou modifiée n'alimente la production sans validation sur un
        environnement de test avec des données représentatives.
\end{itemize}

\begin{table}[H]
\centering
\small
\renewcommand{\arraystretch}{1.5}
\caption{Matrice de contrôle des interfaces inter-applicatives}
\label{tab:controles-interfaces}
\begin{tabular}{|p{2.6cm}|p{3.8cm}|p{3.8cm}|p{2.5cm}|}
    \hline
    \rowcolor{gray!20}
    \textbf{Objectif} & \textbf{Critère} & \textbf{Test d'audit} &
    \textbf{Risque si absent} \\ \hline
    \textbf{Intégrité des flux} &
    Totaux de contrôle calculés côté source et vérifiés côté destination~;
    les écarts déclenchent une alerte &
    Simuler un flux incomplet~; vérifier que l'alerte est émise et reçue
    par le responsable &
    Pertes ou altérations non détectées~; décisions sur données
    incomplètes \\ \hline
    \textbf{Exhaustivité} &
    L'application destination détecte les ruptures de séquence et les
    doublons &
    Soumettre un flux avec un enregistrement manquant~; vérifier la
    détection et l'alerte &
    Enregistrements perdus ou dupliqués intégrés silencieusement \\ \hline
    \textbf{Confidentialité} &
    Les flux sensibles sont chiffrés en transit (TLS~1.2 minimum) &
    Analyser la configuration du protocole~; capturer un échange et
    vérifier qu'il est illisible en clair &
    Interception des données en transit~; violation RGPD \\ \hline
    \textbf{Authentification de l'émetteur} &
    L'identité de l'application source est vérifiée avant acceptation
    du flux &
    Soumettre un flux depuis une source non autorisée~; vérifier le
    rejet &
    Injection de données frauduleuses~; usurpation applicative \\ \hline
    \textbf{Non-répudiation} &
    Un journal horodaté enregistre émetteur, destinataire et résultat
    de chaque transfert &
    Consulter le journal~; vérifier sa complétude et son intégrité &
    Impossibilité de prouver ou contester un transfert en cas de litige
    \\ \hline
    \textbf{Traitement des anomalies} &
    Un responsable désigné surveille les interfaces et traite les
    anomalies dans les délais définis &
    Demander le registre des erreurs (3~derniers mois)~; vérifier que
    chaque anomalie a un plan de correction &
    Anomalies ignorées~; corruptions accumulées sans détection \\ \hline
    \textbf{Procédure de rejeu} &
    En cas d'échec, la relance est possible sans doublon ni perte de
    données &
    Simuler un échec~; rejouer~; vérifier exhaustivité et absence de
    doublons &
    Perte définitive ou duplication silencieuse lors de la reprise \\ \hline
    \textbf{Test avant production} &
    Toute interface nouvelle ou modifiée est testée sur un environnement
    de recette avant mise en production &
    Demander les PV de test des 3~dernières mises en production
    d'interface &
    Anomalies de format ou de logique découvertes sur données réelles
    \\ \hline
\end{tabular}
\end{table}

\begin{boxalert}
\textbf{Point de vigilance~: les interfaces non cartographiées.}

Les interfaces les plus risquées sont celles que la cartographie officielle ne
recense pas~: exports Excel recopiés manuellement, transferts par clé USB,
échanges de données par messagerie. Ces \emph{interfaces informelles} présentent
trois caractéristiques défavorables~:
\begin{enumerate}
    \item \textbf{Aucun contrôle intégré~:} pas de total de contrôle, pas de
        détection de doublon, pas de journal~;
    \item \textbf{Dépendance à la rigueur individuelle~:} la qualité du transfert
        repose entièrement sur l'attention de la personne qui l'exécute~;
    \item \textbf{Invisibilité~:} elles n'apparaissent ni dans la documentation
        applicative, ni dans les contrats de service, ni dans les revues de contrôle
        interne.
\end{enumerate}
Pour les détecter, l'auditeur interroge directement les utilisateurs~:
\emph{«~Comment transmettez-vous les données de cette application à votre collègue
de la comptabilité~?~»} La réponse \og{}on exporte en Excel et on envoie par mail\fg{}
est une réponse \texttt{NON~CONFORME} systématique, quelle que soit la qualité des
interfaces officielles.
\end{boxalert}

\begin{boxlizbiz}
\textbf{Cas Lizbiz's~: l'interface logistique-comptabilité}

L'application de gestion des stocks de Lizbiz's est interfacée avec le module
comptable~: chaque soir, un fichier automatique transfère les mouvements de stock
valorisés vers la comptabilité. L'auditeur vérifie le contrôle d'intégrité~: le
fichier émis contient un total de contrôle, mais l'application comptable ne le
vérifie pas à la réception — elle importe silencieusement ce qu'elle reçoit. Un
test simple le confirme~: l'auditeur modifie manuellement le total dans le fichier
de test~; la comptabilité importe le fichier sans aucune alerte. Par ailleurs,
deux utilisateurs du service logistique avouent envoyer par mail des corrections
de valorisation en dehors du fichier automatique~: interface informelle, non
journalisée, non contrôlée. L'auditeur ne peut pas se prononcer sur la fiabilité
des données comptables tant que ces interfaces ne sont pas corrigées et
documentées.
\end{boxlizbiz}

\subsection{Fréquence et déclencheurs}

La fréquence recommandée est de deux à trois ans pour une application de gestion
stable. Plusieurs événements doivent déclencher un audit anticipé~: une migration
vers une nouvelle version majeure, un changement de prestataire de maintenance,
un incident de sécurité impliquant l'application, une évolution réglementaire
significative, ou une croissance rapide du volume de transactions.

\subsection{Référentiels mobilisés}

\norme{COBIT~2019} fournit le cadre de gouvernance~; \norme{ISO/IEC~27034} traite
spécifiquement de la sécurité applicative~; \norme{ISO~19011} guide la conduite de
la mission d'audit. Pour les applications financières, les normes \norme{IAS/IFRS}
et les exigences de l'auditeur financier externe constituent des référentiels
complémentaires incontournables.

\begin{boxlizbiz}
\textbf{LiZbiZ, Pourquoi auditer l'APPLICATION maintenant~?}

LiZbiZ exploite depuis quatre ans un progiciel de gestion commerciale et comptable
baptisé \textit{APPLICATION}. Deux signaux ont déclenché la mission~: une
discordance entre les totaux de la balance comptable et les relevés bancaires sur
le trimestre écoulé, et la découverte fortuite par un administrateur qu'un
utilisateur avait modifié une facture validée sans qu'aucune trace n'en subsiste
dans les logs. Ces deux constats pointent simultanément une défaillance dans les
contrôles d'application (absence de piste d'audit sur le module comptable) et une
faiblesse dans les contrôles généraux (politique de gestion des droits insuffisante).
L'auditeur doit déterminer si la discordance résulte d'une erreur de traitement,
d'une manipulation intentionnelle ou d'un bug applicatif.
\end{boxlizbiz}


% ==============================================================================
\section{Audit de la Fonction Informatique (DSI)}
\label{sec:audit-dsi}
% ==============================================================================

\subsection{Enjeu fondamental}

La Direction des Systèmes d'Information est le pivot organisationnel de toute
la chaîne de valeur numérique de l'entreprise. Elle conçoit, déploie, exploite
et maintient l'ensemble des actifs informatiques dont dépendent les processus
métiers. Auditer la DSI, ce n'est pas auditer une machine ni un logiciel~: c'est
examiner une \textit{organisation}, sa gouvernance, sa stratégie, ses ressources,
ses méthodes, sa culture et son positionnement au sein de l'entreprise.

L'enjeu central est celui de la \termecle{maturité informatique}~: la DSI est-elle
organisée, dotée, pilotée et régulée de façon à remplir efficacement sa mission
et à créer de la valeur pour l'organisation~? Une DSI immature génère des risques
en cascade~: projets dépassant les délais et les budgets, applications mal maintenues,
incidents récurrents non résolus, dépendance critique vis-à-vis de prestataires
extérieurs ou de personnes-clés dont le départ pourrait paralyser l'organisation.

\subsection{Périmètre et frontières}

Cet audit couvre six domaines~: le rôle et le positionnement de la DSI dans
l'organigramme~; sa planification stratégique et l'existence d'un schéma directeur~;
la maîtrise de ses budgets et de ses coûts~; la mesure de sa performance~; son
organisation interne et la gestion de ses ressources humaines~; enfin, sa conformité
au cadre législatif et réglementaire applicable.

La frontière avec l'audit des applications est claire~: on s'intéresse ici non pas
à ce que produisent les systèmes, mais à la qualité de l'entité qui les produit et
les gère. La frontière avec l'audit de sécurité est plus poreuse~: la politique de
sécurité des SI relève organiquement de la DSI, et sa qualité sera examinée lors
des deux types d'audit, sous l'angle organisationnel dans l'audit de la fonction
informatique, sous l'angle technique dans l'audit de sécurité.

\subsection{Dimensions caractéristiques}

La \termecle{séparation des tâches} est le premier principe structurant. Elle doit
être vérifiée à plusieurs niveaux~: entre les équipes de développement et celles
de la production, entre l'administration des bases de données et l'exploitation
courante, entre la maîtrise d'ouvrage (MOA) et la maîtrise d'œuvre (MOE). Sa
violation est un vecteur classique de fraude et d'erreur non détectée.

La \termecle{dépendance aux personnes-clés} est un risque spécifiquement camerounais
et africain que l'auditeur doit systématiquement évaluer. Lorsqu'une seule personne
détient la connaissance d'un système critique, configuration d'un serveur, paramétrage
d'une application, accès à un prestataire, le départ de cette personne peut paralyser
l'organisation. L'auditeur recherche l'existence de procédures documentées, de
formations croisées et de plans de succession.

Le \termecle{turn-over} de la DSI constitue un indicateur de santé organisationnelle.
Un taux inférieur à 5\,\% par an peut signaler une organisation figée~; supérieur
à 15\,\%, il révèle une instabilité préoccupante susceptible de compromettre la
continuité du service.

La conformité au \textbf{cadre législatif camerounais} est un point de contrôle
distinctif. Les lois de 2010 sur la cybersécurité, la cybercriminalité et les
transactions électroniques imposent des obligations précises que la DSI doit
connaître et faire respecter. L'auditeur vérifiera que ces textes sont connus des
responsables informatiques et que des mesures préventives concrètes ont été mises
en place~: fraude informatique, chiffrement, signature électronique, archivage
électronique.

\subsection{Fréquence et déclencheurs}

La fréquence recommandée est de trois ans. Les déclencheurs d'un audit anticipé
sont~: un changement de DSI ou de directeur général, une restructuration
organisationnelle majeure, une série d'incidents IT récurrents, ou un audit
financier externe ayant relevé des faiblesses dans l'environnement de contrôle
informatique.

\subsection{Référentiels mobilisés}

\norme{COBIT~2019} est la référence principale, particulièrement pour les domaines
de gouvernance (EDM) et de gestion (APO, BAI, DSS, MEA). \norme{ITIL~4} complète
COBIT sur les processus opérationnels. \norme{ISO/IEC~20000-1} traite du management
des services informatiques. \norme{ISO~19011} encadre la conduite de la mission.

\begin{boxlizbiz}
\textbf{LiZbiZ, Une DSI sans gouvernance formalisée}

La DSI de LiZbiZ, composée de sept personnes, est dirigée par un responsable
informatique qui cumule les fonctions de développeur senior, d'administrateur
système et d'interlocuteur contractuel des prestataires. Aucune charte informatique
ne définit le périmètre de ses responsabilités. Le schéma directeur date de 2019
et n'a pas été mis à jour depuis la décision d'intégrer l'IA dans la logistique.
Le budget IT est intégré au budget général sans ligne de dépense détaillée. Un
auditeur n'ayant examiné que l'APPLICATION aurait pu imputer les défaillances
constatées à un problème logiciel~; l'audit de la DSI révèle que leur cause racine
est organisationnelle.
\end{boxlizbiz}


% ==============================================================================
\section{Audit des Projets Informatiques}
\label{sec:audit-projets}
% ==============================================================================

\subsection{Enjeu fondamental}

Les projets informatiques sont des investissements à haut risque. Les études
internationales sur l'échec des projets IT, dont le célèbre \textit{Chaos Report}
du Standish Group, documentent depuis trente ans que moins de 30\,\% des projets IT
sont livrés dans les délais, dans les budgets et avec le périmètre fonctionnel initial.
Les causes de ces échecs sont rarement techniques~: elles sont organisationnelles,
méthodologiques et relationnelles.

L'enjeu de l'audit de projet est donc de s'assurer que les mécanismes de pilotage,
de qualité et de contrôle sont en place pour maximiser les chances de succès, ou,
si le projet est déjà terminé, de documenter les causes d'échec et d'en tirer les
enseignements. La question centrale est~: \textit{ce projet est-il maîtrisé~?}

\subsection{Périmètre et frontières}

L'audit de projet peut intervenir à deux moments distincts. L'\textbf{audit de jalons}
(\textit{gate review}) intervient en cours de projet, à des étapes prédéfinies, pour
vérifier que les conditions de passage à la phase suivante sont réunies~: c'est une
forme d'audit prospectif. L'\textbf{audit de réception} intervient à la clôture pour
vérifier la conformité des livrables aux spécifications et tirer les enseignements
pour les projets futurs. Un projet abandonné ou suspendu doit également être audité.

\subsection{Dimensions caractéristiques}

Le \termecle{triangle QCD} — Qualité, Coût, Délai — est le cadre de référence de
l'audit de projet. L'auditeur examine comment ces trois contraintes ont été définies,
communiquées, suivies et rééquilibrées. La question n'est pas de savoir si les trois
ont été respectées simultanément, mais si les arbitrages ont été conscients, documentés
et validés par les parties prenantes compétentes.

La distinction entre \termecle{MOA} (Maîtrise d'Ouvrage, le client, définisseur des
besoins) et \termecle{MOE} (Maîtrise d'Œuvre, le réalisateur technique) est au cœur
de l'audit de projet. L'auditeur vérifie que ces deux rôles sont clairement séparés,
que leurs responsabilités sont contractuellement définies et que la MOA a effectivement
exercé son rôle de validation à chaque étape.

La \termecle{gestion des risques de projet} constitue un point de contrôle critique.
L'auditeur recherche l'existence d'un registre des risques, régulièrement mis à jour,
avec des plans d'atténuation assignés à des responsables identifiés. L'absence de ce
dispositif est un signal d'alarme fort.

\subsection{Fréquence et déclencheurs}

Tout projet doit être audité à son terme. Les signaux déclencheurs d'un audit anticipé
sont~: un dépassement budgétaire supérieur à 15\,\%, un glissement de planning
supérieur à 20\,\% du délai initial, un changement de chef de projet, ou des conflits
récurrents entre MOA et MOE.

\subsection{Référentiels mobilisés}

\norme{PMBOK} et \norme{PRINCE2} sont les deux référentiels de conduite de projet les
plus répandus. \norme{ISO~21500} propose un cadre normatif international harmonisé.
Pour les projets agiles, \norme{SCRUM} et \norme{SAFe} constituent les cadres de
référence avec des implications spécifiques sur la démarche d'audit, les livrables
n'étant plus définis ex ante mais émergents.

\begin{boxlizbiz}
\textbf{LiZbiZ, Le projet IA sans garde-fous}

LiZbiZ a lancé il y a huit mois un projet d'intégration d'un algorithme
d'apprentissage automatique dans son système de planification logistique. Aucune
charte de projet n'a été formalisée. Le budget alloué a déjà été dépassé de 40\,\%
sans décision formelle de révision. Les données clients utilisées pour entraîner le
modèle n'ont fait l'objet d'aucune analyse de conformité RGPD. La MOA n'a participé
à aucune réunion de validation depuis le lancement. L'audit de jalons révèle un projet
hors de contrôle sur les trois dimensions du triangle QCD, avec un risque juridique
réel lié au traitement non conforme des données personnelles.
\end{boxlizbiz}


% ==============================================================================
\section{Audit du Support Utilisateur et de la Gestion du Parc}
\label{sec:audit-support}
% ==============================================================================

\subsection{Enjeu fondamental}

Le support utilisateur et la gestion du parc informatique constituent le visage
quotidien de la DSI pour l'ensemble des collaborateurs. La qualité de ce service
conditionne directement la productivité des utilisateurs et leur confiance dans les
outils informatiques. Un Helpdesk inefficace, des équipements mal inventoriés, des
licences non maîtrisées ou des procédures de départ lacunaires sont des risques
opérationnels, financiers et de sécurité souvent sous-estimés.

\subsection{Dimensions caractéristiques}

La \termecle{gestion des incidents et des demandes} est évaluée à travers l'existence
et le suivi des \termecle{SLA} (\textit{Service Level Agreements})~: délai de prise
en charge, délai de résolution, taux de résolution au premier contact (\textit{First
Call Resolution Rate}). L'auditeur vérifie que ces indicateurs sont effectivement
mesurés, communiqués aux utilisateurs et que les écarts font l'objet d'un plan
d'amélioration.

La \termecle{procédure d'off-boarding} est un point de contrôle à fort enjeu
sécuritaire. Le départ d'un collaborateur doit déclencher immédiatement la révocation
de l'ensemble de ses accès logiques (réseau, applications, messagerie, VPN) et la
récupération de ses équipements. L'auditeur vérifie l'existence de cette procédure et
son effectivité en testant sur un échantillon de départs récents.

\begin{boxlizbiz}
\textbf{LiZbiZ, Le développeur fantôme}

Trois mois après le départ d'un développeur senior, l'auditeur constate que son
compte VPN est toujours actif et que ses droits d'accès à l'environnement de
production de l'\textit{APPLICATION} n'ont pas été révoqués. Aucune procédure
d'off-boarding formalisée n'existe. Le risque est critique~: cet ancien collaborateur,
qui connaît l'architecture du SI, pourrait se connecter à distance et accéder aux
données clients de LiZbiZ sans qu'aucune alerte ne soit générée.
\end{boxlizbiz}


% ==============================================================================
\section{Audit de Sécurité Informatique}
\label{sec:audit-securite}
% ==============================================================================

\subsection{Enjeu fondamental}

L'audit de sécurité est le type d'audit SI qui a connu la plus forte croissance au
cours des dix dernières années, sous l'effet conjugué de la multiplication des
cyberattaques, de la généralisation de la transformation numérique et du durcissement
des exigences réglementaires. Son enjeu est de mesurer objectivement la résistance
du SI aux menaces, qu'elles soient externes (hackers, groupes criminels organisés,
États) ou internes (employés malveillants, erreurs humaines, sous-traitants
négligents).

Ce qui distingue fondamentalement l'audit de sécurité des autres types d'audit SI
est sa dimension \textit{adversariale}~: l'auditeur ne se contente pas de vérifier
l'existence de contrôles, il cherche à démontrer leur insuffisance en adoptant la
posture d'un attaquant. Les tests d'intrusion (\textit{pentests}) en sont l'expression
la plus radicale.

% ------------------------------------------------------------------------------
\subsection{Les Piliers D.I.C.T.}
\label{subsec:dict}
% ------------------------------------------------------------------------------

L'audit de sécurité s'organise autour des quatre propriétés fondamentales que tout
SI doit garantir.

La \termecle{Disponibilité} vise à s'assurer que le système est accessible aux
utilisateurs autorisés quand ils en ont besoin, y compris dans des conditions
dégradées. Elle est menacée par les pannes matérielles, les saturations de ressources,
les attaques DDoS et les catastrophes naturelles.

L'\termecle{Intégrité} garantit que les données ne peuvent être modifiées, créées
ou supprimées que par des utilisateurs autorisés, selon des processus contrôlés.
Sa violation est parfois difficile à détecter~: une modification subtile d'un montant
comptable peut passer inaperçue pendant des mois si aucun contrôle de cohérence
n'est en place.

La \termecle{Confidentialité} assure que l'information n'est accessible qu'aux
personnes qui en ont le droit. Elle est menacée par des configurations défectueuses,
des droits d'accès excessifs, des données en clair sur des supports non chiffrés ou
des communications non sécurisées.

La \termecle{Traçabilité}, parfois appelée imputabilité, garantit que toute action
sur le SI peut être attribuée à un acteur identifié et horodatée de façon fiable.
Sans elle, il est impossible d'investiguer un incident, de reconstituer une chronologie
d'attaque ou d'apporter la preuve d'une fraude devant une juridiction.

% ------------------------------------------------------------------------------
\subsection{Disponibilité Physique~: les Contrôles Environnementaux}
\label{subsec:controles-environnementaux}
% ------------------------------------------------------------------------------

La disponibilité d'un SI ne dépend pas seulement de son architecture logique,
redondance, failover, PRA, mais aussi des conditions physiques dans lesquelles ses
équipements opèrent. Un centre de calcul parfaitement redondant sur le plan réseau
peut être mis hors service en quelques minutes par une coupure électrique non
anticipée, une inondation ou un incendie. Les contrôles environnementaux constituent
le \emph{socle physique} de la disponibilité~: ils sont la condition préalable à
laquelle tous les contrôles logiques se trouvent subordonnés.

La démarche d'audit environnemental repose sur un principe de
\textbf{défense en profondeur spatiale}~: l'auditeur raisonne par zones concentriques,
du périmètre extérieur jusqu'à la baie critique, en vérifiant que chaque zone dispose
de ses propres dispositifs de protection indépendants des zones adjacentes.

\begin{boxdef}
\textbf{Plan d'Occupation des Sols du SI (POS-SI)~:}
Document cartographiant la localisation physique de l'ensemble des équipements
informatiques au sein des locaux de l'organisation~: salles serveurs, locaux techniques,
postes de travail, équipements réseau, baies de brassage. Il identifie les
\emph{zones fonctionnelles} (production, développement, accès public) et les
\emph{responsables de zones}. L'auditeur utilise ce document comme base de la revue
des contrôles environnementaux~: sans cartographie physique du SI, l'audit ne peut
couvrir que ce que l'audité consent à montrer.
\end{boxdef}

\subsubsection{Les Zones de Sécurité Physique}

L'approche par zones concentriques distingue trois niveaux de protection~:

\begin{enumerate}
    \item \textbf{Zone périphérique (périmètre du bâtiment)~:} contrôle d'accès
        général au bâtiment, gardiennage, vidéosurveillance des entrées, gestion
        des badges visiteurs. La défaillance de cette zone expose l'ensemble du SI
        aux menaces physiques externes.
    \item \textbf{Zone intermédiaire (locaux techniques)~:} accès restreint aux
        personnels habilités, serrures électroniques ou mécaniques sur les locaux
        réseau (baies de brassage, équipements actifs), alarme anti-intrusion. Les
        équipements réseau placés dans des couloirs ou des faux-plafonds accessibles
        constituent une vulnérabilité fréquente.
    \item \textbf{Zone critique (salle serveur)~:} niveau de protection maximal.
        L'accès est nominatif, journalisé et revu périodiquement. Aucune personne
        externe ne doit y accéder seule~; la règle du \emph{double regard}
        (\textit{two-person rule}) est recommandée pour les interventions sensibles.
\end{enumerate}

\paragraph{Clé de contrôle.}
L'auditeur demande le journal d'accès à la salle serveur des 30~derniers jours et
le rapproche de la liste du personnel habilité. Tout accès d'une personne non
habilitée, ou dont l'habilitation a expiré (départ de l'entreprise, changement de
poste), constitue une non-conformité immédiate.

\subsubsection{Les Dispositifs de Protection contre les Menaces Physiques}

L'auditeur réalise une \textbf{visite physique de la salle serveur}. Aucun document
ne remplace l'observation directe~: des détecteurs de fumée peuvent figurer dans
l'inventaire et être non fonctionnels~; une UPS peut être en service mais non
dimensionnée pour la charge réelle.

\begin{table}[H]
\centering
\small
\renewcommand{\arraystretch}{1.5}
\caption{Contrôles environnementaux~: points de vérification lors de la visite
de la salle serveur}
\label{tab:controles-env}
\begin{tabular}{|p{2.8cm}|p{4.8cm}|p{4.5cm}|}
    \hline
    \rowcolor{gray!20}
    \textbf{Menace} & \textbf{Dispositif attendu} &
    \textbf{Test d'audit} \\ \hline
    \textbf{Coupure électrique} &
    ASC/UPS dimensionnée pour couvrir le délai de bascule vers le groupe
    électrogène~; groupe électrogène avec réserve de carburant documentée
    et testée &
    Vérifier le rapport du dernier test de bascule UPS~; vérifier le niveau
    de carburant et la date du dernier essai du groupe \\ \hline
    \textbf{Incendie} &
    Détecteurs de fumée à double technologie~; système d'extinction
    automatique adapté à l'environnement électronique (gaz inerte,
    \textbf{non} eau)~; extincteurs CO\textsubscript{2} à portée &
    Vérifier les rapports de maintenance annuelle des détecteurs et du
    système d'extinction~; vérifier la date de contrôle des extincteurs
    \\ \hline
    \textbf{Eau / inondation} &
    Détecteurs d'humidité au niveau du sol~; absence de tuyauterie
    au-dessus des équipements critiques~; équipements sur plancher surélevé
    ou supports à $\geq 15$~cm du sol &
    Observer le cheminement des tuyauteries~; vérifier la hauteur des
    équipements~; tester les détecteurs d'eau \\ \hline
    \textbf{Chaleur / humidité} &
    Climatisation redondante~; sondes de surveillance avec alertes
    automatiques dans les plages recommandées par les constructeurs &
    Demander le journal de température et d'hygrométrie des 30~derniers
    jours~; vérifier l'existence d'une alarme en cas de dépassement \\ \hline
    \textbf{Électricité statique} &
    Revêtement de sol antistatique ou dalles antistatiques~; bracelets
    antistatiques disponibles pour les interventions matérielles &
    Observation directe lors de la visite \\ \hline
    \textbf{Accès non autorisé} &
    Contrôle d'accès électronique nominatif~; journal d'accès horodaté~;
    caméras de surveillance à l'entrée de la salle &
    Comparer le journal d'accès à la liste des personnes habilitées~;
    vérifier que les comptes révoqués ont bien été désactivés \\ \hline
\end{tabular}
\end{table}

\subsubsection{La Destruction Sécurisée des Supports}

La fin de vie d'un équipement ou d'un support de stockage constitue un vecteur de
fuite d'informations souvent sous-estimé. Un disque dur réformé et revendu sans
effacement préalable peut contenir des données confidentielles, des identifiants ou
des clés cryptographiques. L'auditeur vérifie que~:

\begin{itemize}
    \item Une procédure formalisée encadre la mise au rebut de tout équipement
        contenant des données~: ordinateurs, serveurs, disques de sauvegarde,
        imprimantes multifonctions (dont le disque interne est souvent ignoré),
        smartphones professionnels~;
    \item L'effacement est réalisé selon une méthode certifiée adaptée au support
        (écrasement multiple conforme à la norme DoD~5220.22-M pour les disques
        magnétiques~; dégaussage ou destruction physique pour les supports très
        sensibles~; effacement cryptographique pour les SSD)~;
    \item Des \textbf{certificats de destruction} sont émis et conservés pour chaque
        opération, constituant une preuve opposable en cas de litige ou d'audit
        réglementaire~;
    \item Les documents papier confidentiels sont détruits par broyage certifié
        niveau DIN~66399-P4 minimum (particules $\leq 160~\text{mm}^2$), et non par
        simple déchiquetage en bandelettes.
\end{itemize}

\begin{boiteRemarque}{Le cas particulier des supports cloud et virtualisés}
Dans un environnement virtualisé ou cloud, la destruction physique des supports
n'est pas possible. La parade consiste en l'\textbf{effacement cryptographique}~:
les données sont chiffrées dès leur création, et la destruction de la clé les rend
inaccessibles de façon irréversible. L'auditeur vérifie que la politique de
chiffrement dès la création (\textit{encrypt-at-rest by default}) est appliquée et
que la gestion des clés est documentée et auditée.
\end{boiteRemarque}

\begin{boxlizbiz}
\textbf{Cas Lizbiz's~: la salle serveur sans plancher surélevé}

Lors de la visite physique, l'auditeur constate que la salle serveur de Lizbiz's
est située en rez-de-chaussée, les équipements directement posés sur le sol carrelé,
sans plancher surélevé. Le bâtiment est situé en zone urbaine dense, et les archives
de la ville font état de deux épisodes de remontée des eaux dans le quartier au cours
des dix dernières années. Le plan de continuité prévoit un RTO de 4~heures~; or,
une inondation même mineure de 5~cm suffirait à mettre hors service l'ensemble des
équipements. Ce constat, obtenu en dix minutes de visite, invalide la crédibilité du
PCA entier~: le risque physique n'avait pas été intégré dans le BIA.
Statut~: \texttt{NON~CONFORME}. Recommandation~: surélévation des baies à $\geq
15$~cm ou relocalisation dans un local non exposé, délai~\texttt{IMMÉDIAT}.
\end{boxlizbiz}

% ------------------------------------------------------------------------------
\subsection{Dimensions Caractéristiques de l'Audit de Sécurité}
\label{subsec:securite-dimensions}
% ------------------------------------------------------------------------------

L'audit de sécurité mobilise trois méthodologies complémentaires dont la combinaison
dépend du périmètre et des objectifs de la mission.

L'\textbf{audit de configuration} consiste à vérifier que les paramètres des systèmes
(serveurs, équipements réseau, applications) sont conformes aux référentiels de
durcissement (\textit{hardening}) reconnus, tels que les benchmarks \norme{CIS}
(\textit{Center for Internet Security}). Il détecte les mots de passe par défaut
non modifiés, les services inutiles exposés, les correctifs de sécurité non appliqués.

Les \textbf{tests d'intrusion} (\textit{pentests}) simulent une attaque réelle dans un
cadre contractuellement défini~: périmètre, règles d'engagement, \textbf{autorisation
écrite préalable obligatoire}. Ils peuvent être menés en boîte noire (l'auditeur ne
connaît pas l'infrastructure), en boîte grise (informations partielles) ou en boîte
blanche (accès complet à la documentation). Le pentest ne remplace pas l'audit de
configuration~; il le complète en démontrant l'exploitabilité réelle des vulnérabilités
identifiées.

L'\textbf{audit physique} examine la sécurité des locaux~: contrôle des accès à la
salle des serveurs, gestion des visiteurs, destruction des documents confidentiels,
protection contre les sinistres. Il est souvent négligé mais constitue régulièrement
le maillon faible de la chaîne de sécurité — comme illustré par les contrôles
environnementaux présentés à la section précédente.

\subsection{Fréquence et déclencheurs}

La fréquence minimale recommandée est annuelle. Les tests de pénétration externes
sur les périmètres exposés à Internet peuvent être plus fréquents, voire continus
dans les organisations les plus matures (programme de \textit{bug bounty}). Les
déclencheurs d'un audit non planifié sont~: tout incident de sécurité significatif,
une modification majeure de l'architecture, l'ouverture d'un nouveau périmètre
réseau, l'intégration d'un prestataire tiers.

\subsection{Référentiels mobilisés}

\norme{ISO/IEC~27001:2022} est le référentiel de management de la sécurité de
l'information. \norme{NIST Cybersecurity Framework~v2.0} (2024) propose un cadre
fonctionnel articulé autour de six fonctions~: Gouverner, Identifier, Protéger,
Détecter, Répondre, Rétablir. \norme{ISO/IEC~27035} traite de la gestion des
incidents. Pour les tests d'intrusion, le référentiel \norme{PTES} et
l'environnement \norme{OWASP} pour les applications web font autorité.

\begin{boxlizbiz}
\textbf{LiZbiZ, Trois observations critiques}

Le pentest mené sur le périmètre de LiZbiZ révèle~: \textit{(i)}~le serveur Apache
hébergeant l'\textit{APPLICATION} présente une vulnérabilité critique non corrigée
depuis six mois, permettant une exécution de code à distance~; \textit{(ii)}~le
serveur de base de données comptable est sur le même segment réseau que les postes
utilisateurs, supprimant toute barrière entre un poste compromis et les données
financières~; \textit{(iii)}~les journaux du pare-feu ne sont conservés que sept jours,
rendant impossible toute investigation rétrospective sur un incident antérieur.
Ces trois constats sont interdépendants~: un attaquant ayant exploité la vulnérabilité
Apache peut atteindre latéralement la base de données sans que personne ne le sache
ni ne puisse le prouver.
\end{boxlizbiz}


% ==============================================================================
\section{Audit de la Continuité d'Activité et de la Résilience}
\label{sec:audit-continuite}
% ==============================================================================

\subsection{Enjeu fondamental}

La résilience organisationnelle désigne la capacité d'une entité à anticiper les
disruptions, à y résister et à reprendre ses activités critiques dans des délais
compatibles avec sa survie. L'audit de la continuité d'activité ne se demande pas
\textit{si} un sinistre surviendra~; il postule qu'il surviendra, et se demande
\textit{si l'organisation est prête à y faire face}.

Cet audit est en forte croissance depuis les cyberattaques par rançongiciel qui ont
paralysé des hôpitaux, des administrations et des entreprises entières. Il est
également requis par un nombre croissant de régulateurs sectoriels qui imposent des
exigences formelles de continuité.

\subsection{Concepts Fondamentaux}

Quatre indicateurs structurent l'évaluation de la résilience.

Le \termecle{BIA} (\textit{Business Impact Analysis}) est le point de départ~: il
identifie les processus critiques, quantifie l'impact d'une interruption sur chacun
d'eux et détermine les ressources minimales nécessaires à leur maintien. L'auditeur
vérifie l'existence, la complétude et la mise à jour régulière du BIA.

Le \termecle{MTPD} (\textit{Maximum Tolerable Period of Disruption}) est la durée
maximale d'interruption d'un processus au-delà de laquelle la survie de l'organisation
est menacée. C'est une donnée métier, déterminée par les dirigeants, non par la DSI.

Le \termecle{RTO} (\textit{Recovery Time Objective}) est l'objectif de délai de
remise en service après un sinistre. Il doit être inférieur au MTPD~: un RTO supérieur
au MTPD est une incohérence fondamentale que l'auditeur doit impérativement relever.

Le \termecle{RPO} (\textit{Recovery Point Objective}) est la perte de données
maximale admissible, exprimée en unité de temps. Il est directement lié à la fréquence
des sauvegardes~: un RPO de quatre heures implique des sauvegardes au moins toutes les
quatre heures.

\begin{boxalert}
\textbf{L'incohérence RTO/RPO~: un classique de l'audit.}
Un RTO de 4~h avec un RPO de 24~h est techniquement incohérent dans de nombreux
contextes~: on peut relancer les systèmes en quatre heures, mais avec des données
vieilles d'une journée, ce qui peut être inacceptable pour un processus financier ou
logistique critique. L'auditeur doit systématiquement vérifier la cohérence de ces
paramètres avec les exigences métiers réelles.
\end{boxalert}

\subsection{Fréquence et déclencheurs}

La fréquence recommandée est annuelle, incluant obligatoirement un \textbf{test réel}
du PCA et du PRA~: un plan non testé est un plan dont on ignore s'il fonctionnera.
Les déclencheurs d'un audit non planifié sont~: tout incident ayant nécessité
l'activation du PCA/PRA, une modification significative de l'architecture du SI,
ou une nouvelle exigence réglementaire.

\subsection{Référentiels mobilisés}

\norme{ISO~22301:2019} est le standard international du management de la continuité
d'activité. \norme{ISO~22317} traite spécifiquement du BIA. \norme{NIST SP~800-34}
offre un guide pratique de planification de la continuité pour les systèmes
d'information.

\begin{boxlizbiz}
\textbf{LiZbiZ, Un PCA sur papier, jamais testé}

LiZbiZ dispose d'un Plan de Continuité d'Activité rédigé en~2021. Il identifie
correctement les processus critiques et définit un RTO de 4~h et un RPO de 8~h pour
l'\textit{APPLICATION}. Cependant, l'auditeur constate~: \textit{(i)}~le PCA n'a
jamais été testé~; \textit{(ii)}~les sauvegardes sont effectuées quotidiennement à
23~h, générant un RPO réel de 23~h, incompatible avec l'objectif de 8~h~;
\textit{(iii)}~les coordonnées de l'hébergeur de secours correspondent à un
prestataire dont le contrat a été résilié en~2022. Le PCA de LiZbiZ est un document
de conformité apparente, pas un dispositif opérationnel.
\end{boxlizbiz}


% ==============================================================================
\section{Audit de la Fonction Étude (Développement)}
\label{sec:audit-etude}
% ==============================================================================

\subsection{Enjeu fondamental}

La fonction étude est responsable de la conception, du développement et de la
maintenance évolutive des applications. Son audit examine la rigueur des méthodes
employées, la qualité du code produit, la maîtrise de la documentation et
l'intégration de la sécurité dès la conception. C'est un audit qui requiert à la
fois des compétences organisationnelles et une capacité à apprécier la qualité
technique des livrables.

L'enjeu principal est de s'assurer que les applications produites en interne, ou
développées en sous-traitance, sont fiables, maintenables, documentées et conformes
aux exigences de sécurité. Une fonction étude mal organisée produit une
\textit{dette technique} qui s'accumule silencieusement~: le code non documenté
devient inconnaissable au départ des développeurs, les applications non testées
accumulent des bugs latents.

\subsection{Dimensions caractéristiques}

La \termecle{séparation des environnements} est un contrôle fondamental~: les
environnements de développement, de test et de production doivent être strictement
séparés. Un développeur ne doit jamais avoir accès direct à la production~; toute
mise en production doit passer par une procédure formalisée avec validation
indépendante.

La \termecle{gestion des versions} (\textit{versioning}) assure la traçabilité de
l'évolution du code source. L'auditeur vérifie l'utilisation d'un système de contrôle
de version (Git ou équivalent), la politique de branches, et la capacité à reproduire
un état antérieur du système en cas d'incident.

Le \termecle{Security by Design} — intégration de la sécurité dès la conception —
est le paradigme moderne du développement sécurisé. L'auditeur vérifie l'existence
de normes de codage sécurisé, la réalisation d'analyses statiques du code (SAST)
pour détecter les vulnérabilités classiques (injections SQL, XSS, gestion incorrecte
des erreurs) et l'intégration de tests de sécurité dans le pipeline de déploiement
(DevSecOps).

\subsection{Fréquence et déclencheurs}

La fréquence recommandée est de deux à trois ans, ou lors de l'adoption d'une
nouvelle méthodologie de développement (passage à l'Agile, adoption du DevOps),
de l'intégration d'un nouveau prestataire, ou d'un incident applicatif grave ayant
pour origine une défaillance du processus de développement.

\subsection{Référentiels mobilisés}

\norme{ISO/IEC~25010} définit les caractéristiques de qualité d'un produit logiciel.
\norme{ISO/IEC~27034} traite de la sécurité applicative. \norme{OWASP SAMM}
(\textit{Software Assurance Maturity Model}) propose un modèle de maturité pour les
pratiques de développement sécurisé. Pour les projets agiles, \norme{SAFe} et les
pratiques DevSecOps constituent les références contemporaines.


% ==============================================================================
\section{Audit des Tiers, du Cloud et de l'Externalisation}
\label{sec:audit-tiers}
% ==============================================================================

\subsection{Enjeu fondamental}

L'externalisation informatique — hébergement cloud, infogérance, sous-traitance de
développement, SaaS — transfère des activités et des données à des tiers, mais ne
transfère pas la responsabilité. L'organisation reste juridiquement responsable de
la sécurité et de la conformité de ses données, y compris lorsqu'elles sont traitées
par un prestataire. C'est le principe fondateur du RGPD en matière de sous-traitance.

L'audit des tiers répond à une question simple mais dont la réponse est souvent
inconfortable~: \textit{faites-vous confiance à votre prestataire, ou en avez-vous
la preuve~?}

\subsection{Dimensions caractéristiques}

Le \termecle{droit d'audit} est la première vérification contractuelle. Toute
convention avec un prestataire hébergeant des données sensibles doit contenir une
clause permettant à l'organisation, ou à ses mandataires, d'auditer le prestataire.
Son absence prive l'organisation de tout moyen de contrôle indépendant.

Les \termecle{rapports SOC} (\textit{Service Organization Controls}) constituent le
substitut standard au droit d'audit direct lorsque celui-ci est contractuellement ou
pratiquement impossible, notamment vis-à-vis des grands hyperscalers (AWS, Azure,
GCP). Deux niveaux sont pertinents pour l'auditeur SI~:
\begin{itemize}
    \item \textbf{SOC~1 Type~II~:} évalue les contrôles relatifs au reporting
        financier sur une période d'observation (six ou douze mois). Pertinent pour
        les prestataires hébergeant des applications comptables ou financières.
    \item \textbf{SOC~2 Type~II~:} évalue les contrôles relatifs à la sécurité,
        la disponibilité, l'intégrité du traitement, la confidentialité et la
        protection des données personnelles. C'est le rapport le plus fréquemment
        demandé pour les prestataires cloud hébergeant des données sensibles.
\end{itemize}

L'auditeur ne se contente pas de vérifier l'existence de ces rapports~: il en analyse
le contenu, notamment les exceptions relevées par l'auditeur du prestataire.

La \termecle{cartographie des tiers} est un préalable indispensable~: l'organisation
doit disposer d'un inventaire à jour de tous ses prestataires informatiques, avec pour
chacun la nature des données traitées, le niveau de criticité et les mécanismes de
contrôle en place.

\subsection{Fréquence et déclencheurs}

La revue des tiers doit être annuelle au minimum. Tout changement significatif chez
un prestataire critique (rachat, incident de sécurité public, modification des
conditions contractuelles) doit déclencher une revue immédiate.

\subsection{Référentiels mobilisés}

\norme{ISO/IEC~27036} traite spécifiquement de la sécurité des relations avec les
fournisseurs. \norme{ISO/IEC~27017} et \norme{27018} complètent ISO~27001 pour les
environnements cloud. Le \norme{RGPD} (articles~28 et~29) encadre contractuellement
les relations entre responsables de traitement et sous-traitants.

\begin{boxlizbiz}
\textbf{LiZbiZ, Un cloud sans filet}

LiZbiZ héberge ses données clients et ses sauvegardes sur une solution cloud d'un
prestataire local. L'auditeur constate~: \textit{(i)}~le contrat ne contient aucune
clause d'audit~; \textit{(ii)}~le prestataire ne publie aucun rapport SOC et ne dispose
d'aucune certification ISO~27001~; \textit{(iii)}~aucun inventaire des prestataires
n'existe dans la DSI — le responsable informatique ne peut pas lister l'ensemble des
services tiers utilisés. LiZbiZ fait confiance à son prestataire sans en avoir la
moindre preuve formelle, en violation de ses obligations en tant que responsable de
traitement de données personnelles.
\end{boxlizbiz}


% ==============================================================================
\section{Interactions entre Types d'Audit et Cartographie du Risque SI}
\label{sec:interactions-types}
% ==============================================================================

Les huit types d'audit présentés dans ce chapitre ne sont pas des silos indépendants.
Ils forment un système d'observation complémentaire dont les résultats se nourrissent
mutuellement. Comprendre ces interactions est ce qui distingue l'auditeur expérimenté
du technicien compétent.

L'audit de la fonction informatique fournit le cadre organisationnel dans lequel tous
les autres audits s'inscrivent~: une DSI mal gouvernée génère mécaniquement des
faiblesses dans la gestion des applications, des projets et de la sécurité. Il est
logique de commencer par cet audit lorsqu'une organisation est auditée pour la
première fois.

L'audit de sécurité révèle fréquemment des vulnérabilités dont la cause racine est
organisationnelle~: des droits d'accès excessifs consécutifs à l'absence de procédure
d'off-boarding, des correctifs non appliqués faute de processus de gestion des patches,
une segmentation réseau déficiente résultant d'une architecture non documentée. Ces
causes racines relèvent de l'audit de la fonction informatique, pas de l'audit de
sécurité.

L'audit de continuité complète l'audit de sécurité en posant la question
\textit{après} l'incident~: une fois la cyberattaque survenue, l'organisation est-elle
capable de reprendre ses activités~? Les deux audits sont complémentaires et doivent
idéalement être conduits de façon coordonnée. L'audit des contrôles environnementaux,
présenté dans la section sécurité, nourrit directement l'évaluation du BIA~: un risque
physique non intégré dans le BIA rend le PCA inopérant, comme l'illustre le cas de la
salle serveur de Lizbiz's.

L'audit des applications en service est inséparable de l'audit des interfaces
inter-applicatives~: une application fiable dont les données d'entrée arrivent par des
interfaces non contrôlées ne produit que des résultats dont la fiabilité est illusoire.
L'intégrité bout en bout d'un processus métier dépend de la qualité de l'ensemble des
maillons applicatifs et de leurs connexions.

L'audit des tiers conditionne la portée réelle de tous les autres~: si les données
critiques sont hébergées chez un tiers non audité, les contrôles internes de
l'organisation ne valent que pour la fraction du SI qu'elle maîtrise directement.

\begin{boxresume}
\textbf{À retenir.} Les huit types d'audit SI couvrent un continuum allant de la
gouvernance (audit de la DSI) au fonctionnement (applications, projets, support)
en passant par la sécurité et la résilience (sécurité, continuité) jusqu'à
l'écosystème étendu (tiers et cloud). Chaque type répond à une question fondamentale
distincte, mobilise des référentiels spécifiques et requiert des compétences
particulières. Deux domaines transversaux traversent plusieurs types d'audit~: les
contrôles d'interfaces inter-applicatives, qui conditionnent la fiabilité des données
quel que soit l'applicatif concerné, et les contrôles environnementaux, qui
conditionnent la disponibilité physique de l'ensemble du SI. Le chapitre suivant
décompose chacun de ces types selon le modèle opérationnel du cours, de la définition
du domaine jusqu'à la criticité des conséquences, fournissant ainsi à l'auditeur les
instruments de mesure qui transforment ces typologies en missions concrètes.
\end{boxresume}
